% !TeX encoding = windows-1251
\documentclass[a4paper,12pt]{article}
\usepackage{newlistok}
\usepackage{tikz}

\УвеличитьВысоту{2.3cm}
\УвеличитьШирину{1.5cm}
\renewcommand{\spacer}{\vspace{1mm}}

\Заголовок{Домашняя работа}
\def\словоЛисток{ДЗ \No\/}
\НомерЛистка{3}
\ДатаЛистка{27 февраля 2013г.}

\begin{document}

%\ncopy{1}{
\vspace*{-17mm}
\СоздатьЗаголовок

{\footnotesize
Существует два основных метода решения задач: метод \emph{пристального разглядывания} и метод \emph{тщательного переписывания}.
Используя первый, можно \emph{увидеть} корень уравнения, равные углы, аналогии с другими задачами, решения в одну строчку, решения через другую область математики и т.д.
А во втором вам предстоит потренироваться ;-)
\par}

\задача
На полке в ряд стоят четыре книги. Сколькими способами их можно переставить так, чтобы ни одна не осталась стоять на том же месте?
\кзадача

\задача
На карусели установлено 7 двухместных верблюдов. Сколькими способами можно рассадить на неё 14 школьников?
\кзадача

\задача
Сколькими способами можно представить число 7 в виде суммы нескольких натуральных слагаемых (порядок слагаемых значения не имеет)?
\кзадача

\задача
На чёртовом колесе 11 одинаковых кабинок. 
Сколькими способами можно 7 из них перекрасить в жёлтый цвет?
\кзадача


\hrl
{\small Для получения оценки $n$ необходимо правильно решить $n-1$ задачу.}
%}

\bigskip
\noindent%
\onlyput{-\oddsidemargin}{-300mm}{%
\begin{tikzpicture}[scale=.5]
\draw[step=1,opacity=.5,color=gray] grid (42,60);
\end{tikzpicture}}%
\newpage
\vspace*{-\topmargin}\vspace*{-22mm}%
\noindent%
\onlyput{-\oddsidemargin}{-300mm}{%
\begin{tikzpicture}[scale=.5]
\draw[step=1,opacity=.5,color=gray] grid (42,60);
\end{tikzpicture}}%


\end{document}
